\section{Choice principles used later}
\label{sec:set-prerequisites-06}

\begin{definition}[Axiom of Choice]
The \emph{Axiom of Choice} states that for every set \(X\) of nonempty sets
there exists a function \(c:X\to\bigcup X\) such that \(c(A)\in A\) for every
\(A\in X\).
\end{definition}

\begin{theorem}[Well-ordering theorem]
Assuming the Axiom of Choice, every set can be well-ordered.
\textit{Status: complete proof, \cref{thm:p01-w02-p01-008-well-ordering}.}
\end{theorem}

\begin{corollary}[Cardinal assignment]
\label{cor:cardinal-assignment}
Assuming the Axiom of Choice, every set \(A\) is bijective with a unique
cardinal.
\textit{Status: complete proof, \cref{thm:p01-w02-p01-008-cardinal-assignment}.}
\end{corollary}

Complete proofs are given in \cref{sec:p01-w02-p01-008-choice}.  The
cardinal bijective with \(A\) is denoted by \(\abs A\), and its order agrees with the
injection-based comparison defined above.  The cardinal of \(\N\) is
denoted by \(\aleph_0\).  For a nonempty finite list of
cardinals, \(\max(\kappa_1,\ldots,\kappa_n)\) denotes its greatest member in
the cardinal order.  The well-ordering theorem makes that order total under
the stated Axiom of Choice, so this maximum exists.  These conventions are
the exact cardinal conventions used by the Henkin and
L\"owenheim--Skolem constructions.

\begin{theorem}[Finite coding over an infinite cardinal]
\label{thm:infinite-cardinal-finite-coding}
Assume the Axiom of Choice, and let \(\kappa\) be an infinite cardinal.
Then:
\begin{enumerate}
  \item \(\abs{\kappa^n}=\kappa\) for every positive \(n\in\N\);
  \item the set of all finite sequences from \(\kappa\) has cardinality
        \(\kappa\);
  \item if \(\alpha<\kappa\) and \(A_\beta\) is finite for every
        \(\beta<\alpha\), then
        \(\abs{\bigcup_{\beta<\alpha}A_\beta}<\kappa\); and
  \item if \(\abs{A_n}\leq\kappa\) for every \(n\in\N\), then
        \(\abs{\bigcup_{n\in\N}A_n}\leq\kappa\).
\end{enumerate}
\textit{Status: complete proof, \cref{thm:p01-w02-p01-011-finite-coding}.}
\end{theorem}

\begin{corollary}[Syntax counting]
\label{cor:syntax-counting}
Assume the Axiom of Choice and let \(\kappa\) be an infinite cardinal.
\begin{enumerate}
  \item Every set \(A\) with \(\abs{A}\leq\kappa\) satisfies
        \(\abs{A^{<\omega}}\leq\kappa\), where \(A^{<\omega}\) is the set of
        finite sequences from \(A\).
  \item If \(\Sigma\) is an alphabet with \(\abs{\Sigma}\leq\kappa\), then
        the finite strings over \(\Sigma\), the finite tuples of such
        strings, and the formulas generated by any fixed finitary grammar
        whose alphabet extends \(\Sigma\) by at most countably many further
        symbols each form a set of cardinality at most \(\kappa\).
\end{enumerate}
\textit{Status: complete proof.}
\end{corollary}

\begin{proof}
For item~1, fix an injection \(f:A\to\kappa\); such an injection is the
meaning of \(\abs{A}\leq\kappa\).  The map
\((a_1,\ldots,a_n)\mapsto(f(a_1),\ldots,f(a_n))\) is a length-preserving
injection of \(A^{<\omega}\) into \(\kappa^{<\omega}\), and item~2 of
\cref{thm:infinite-cardinal-finite-coding} gives
\(\abs{\kappa^{<\omega}}=\kappa\).

For item~2, the finite strings over \(\Sigma\) lie in \(\Sigma^{<\omega}\):
apply item~1 with \(A=\Sigma\).  The finite tuples of such strings lie in
\((\Sigma^{<\omega})^{<\omega}\): apply item~1 with \(A=\Sigma^{<\omega}\),
a set of cardinality at most \(\kappa\) by the previous step.  Finally,
every expression generated by a finitary grammar is a finite string over the
grammar's alphabet (\cref{thm:p01-w02-p01-011-grammar-strings}).
Write \(X\) for the at most countable set of adjoined symbols hypothesized
in the statement---logical symbols, punctuation, and variables.  Item~4 of
\cref{thm:infinite-cardinal-finite-coding}, applied to the sequence
\(\Sigma,X,\varnothing,\varnothing,\ldots\), gives
\(\abs{\Sigma\cup X}\leq\kappa\), so the formula set injects into
\((\Sigma\cup X)^{<\omega}\) and item~1 applies.
\end{proof}

See \cite[Theorem~3.5 and Chapter~5]{Jech2003} for complete proofs of
\cref{thm:infinite-cardinal-finite-coding}: Theorem~3.5 gives
\(\aleph_\alpha\cdot\aleph_\alpha=\aleph_\alpha\), and with the
well-ordering theorem cited above every infinite cardinal is an aleph;
Chapter~5 gives the cardinal sums and the union bound
\(\abs{\bigcup_{i\in I}A_i}\leq\abs{I}\cdot\sup_{i\in I}\abs{A_i}\) under
the Axiom of Choice.  \Cref{thm:infinite-cardinal-finite-coding}
and \cref{cor:syntax-counting} supply exactly the cardinal arithmetic
invoked by the Henkin and L\"owenheim--Skolem proofs; no claim about
arbitrary infinite products is being made.  Complete proofs that use only the definitions
of this chapter and the Axiom of Choice are given in \cref{sec:p01-w02-p01-011-finite-coding} and
\cref{sec:p01-w02-p01-011-syntax-counting}; there \(\abs X\leq\kappa\) for
the adjoined symbols is derived from \(\N\subseteq\kappa\).

